% !TEX root = ../../main.tex
\subsection{对比分析}

本节将 FoodShield 与现有外卖平台通信方案、普通 WebSocket 通信系统以及学术匿名通信方案进行安全特性对比，如表~\ref{tab:compare}所示。

\begin{table}[H]
\centering
\caption{FoodShield 与现有方案安全特性对比}
\label{tab:compare}
\footnotesize
\setlength{\tabcolsep}{3pt}
\begin{tabular}{lcccc}
\hline
\textbf{安全属性} & \textbf{FoodShield} & \textbf{美团/饿了么} & \textbf{普通WebSocket} & \textbf{群签名方案} \\
\hline
G1 匿名性 & PID强匿名 & 虚拟号码（弱） & 无匿名 & 签名匿名 \\
G2 认证性 & HMAC-SM3 & 无凭证认证 & 无认证 & 群签名认证 \\
G3 机密性 & SM4-CBC加密 & 传输加密+明文 & 传输加密+明文 & 加密（可选） \\
G4 完整性 & SM3+Merkle & 无验证 & 无验证 & 无Merkle \\
G5 可溯源 & 条件溯源 & 实名查询 & 不可溯源 & 群管理打开 \\
G6 抗滥用 & 验证码+限流 & 无机制 & 无机制 & 群成员管理 \\
G7 隐私最小化 & hash-only & 全量可见 & 全量可见 & 条件打开 \\
系统原型 & 已实现 & 商业运行 & 通用实现 & 仅理论 \\
\hline
\end{tabular}
\end{table}

\textbf{与美团/饿了么对比}：现有平台采用虚拟号码和脱敏显示，但存在三个不足：虚拟号码仅保护手机号，骑手仍可通过订单备注和地址推断用户习惯；管理员可无条件查看通信明文，缺乏权限分离；通信记录明文存储，无完整性验证。FoodShield 通过 PID（跨订单不可链接）、SM4 加密存储、hash-only 审计和 Merkle 完整性验证系统性改进了上述问题。

\textbf{与普通 WebSocket 对比}：普通 WebSocket 仅依赖传输层加密，攻击者获取订单号即可进入通信通道（无 Token 认证），管理员可查看全部明文，日志可被任意篡改。FoodShield 在通信层之上叠加了 Token 认证、PID 匿名、SM4 加密和 Merkle 验证，将普通即时通信升级为可认证、可审计的安全通信。

\textbf{与群签名方案对比}：群签名在理论上提供更强的匿名性（不可链接性），但存在三个实际限制：签名/验证涉及多次双线性映射运算，单次耗时数十毫秒，不适合高频低延迟的外卖通信；缺乏 Merkle 等完整性保护机制；未见针对外卖场景的系统原型。FoodShield 虽匿名性理论强度稍逊，但在计算效率、完整性和系统可用性方面更具优势。

\textbf{综合评价}：FoodShield 的核心差异化在于"可控匿名"模型——日常通信中对骑手和攻击者实现 PID 强匿名，同时在投诉驱动和审批约束下保留条件溯源能力。该模型介于"完全实名不可匿"与"完全匿名不可治"之间，在外卖订单通信场景中实现了匿名性、认证性、机密性、完整性、可溯源性和隐私最小化的多目标平衡。
